\documentclass[11pt]{article}
\usepackage{amssymb,amsmath,amsthm}
\usepackage[urlcolor=blue,colorlinks=true]{hyperref}
\usepackage{mathtools}

\newtheorem{theorem}{Theorem}
\newtheorem{lemma}[theorem]{Lemma}
\newtheorem{corollary}[theorem]{Corollary}
\newtheorem{conjecture}[theorem]{Conjecture}
\newtheorem{observation}[theorem]{Observation}
\theoremstyle{definition}
\newtheorem{definition}[theorem]{Definition}
\theoremstyle{remark}
\newtheorem{remark}[theorem]{Remark}

\newcommand{\cF}{\mathcal{F}}
\newcommand{\cP}{\mathcal{P}}
\newcommand{\emp}{\varnothing}
\newcommand{\Nhalf}{\tfrac{1}{2}}

\title{Frankl's Union-Closed Sets Conjecture:\\
Structural Bounds and Finite Verification}
\author{Independent research project}
\date{\today}

\begin{document}

\maketitle

\begin{abstract}
We study Frankl's union-closed sets conjecture (1979): every finite
union-closed family of sets, other than $\{\varnothing\}$, contains an element
belonging to at least half of its member-sets. The conjecture remains open;
the best unconditional general lower bound is a constant
$c\approx 0.3823455$ (Cambie 2025), and the entropy-based method leading to
it is provably sharp at that constant, so a new idea is required to reach
$1/2$.

We present three elementary structural results, each with a complete proof:
(1) if the family contains a singleton $\{x\}$, then $x$ is popular
(density $\ge \Nhalf$); (2) an injection bound relating the subfamilies of
sets containing a member-set $S$ and sets disjoint from $S$; and (3) a degree
bound $M\ge m/(k+1)$ where $m$ is the number of member-sets and $k$ the size
of a smallest non-empty member-set. We also give exact finite verification
for all union-closed families on at most five ground elements and discuss the
precise obstruction to the full conjecture.

\textbf{Status:} the full conjecture is unresolved. No new theorem resolving
Frankl's conjecture is claimed. The results herein are rigorous partial/
structural results, clearly distinguished from the open conjecture.
\end{abstract}

\section{Introduction}

A family $\cF$ of finite sets is \emph{union-closed} if
$A,B\in\cF\Rightarrow A\cup B\in\cF$. In 1979, P.~Frankl conjectured that
every finite union-closed family $\cF\neq\{\varnothing\}$ has an element
contained in at least half of its member-sets:

\begin{conjecture}[Frankl 1979]
\label{conj:frankl}
If $\cF\subseteq\cP([n])$ is union-closed and $\cF\neq\{\varnothing\}$, then
there exists $x\in[n]$ such that
$$\left|\{A\in\cF:x\in A\}\right|\ge \frac{|\cF|}{2}.$$
\end{conjecture}

An element satisfying the conclusion is called \emph{popular} (or
\emph{abundant}). Despite its elementary statement, the conjecture has
resisted proof for over forty years and is one of the best-known open problems
in extremal set theory.

\subsection*{Why the problem is hard}

Write $d(x)=|\{A\in\cF:x\in A\}|$ and $m=|\cF|$. Since
$\sum_x d(x)=\sum_{A\in\cF}|A|=m\bar s$ where $\bar s$ is the average
member-set size, if $\bar s\ge n/2$ then the average density of an element is
at least $\Nhalf$, so some element is popular. Thus a counterexample must have
$\bar s<n/2$: its member-sets are on average small, yet some element must still
appear in at least half the sets. This tension is the core difficulty.

The strongest unconditional result is a constant lower bound
$c\approx0.3823455$ (Cambie~\cite{cambie2025}, improving
$(3-\sqrt5)/2\approx0.381966$ of Sawin~\cite{sawin2022}, itself building on
Gilmer's first constant bound~\cite{gilmer2022}). Cambie shows the entropy
method behind these bounds is \emph{provably sharp} at
$c\approx0.3823455$: a matching extremal distribution defeats any strict
improvement within that framework. Hence reaching $1/2$ requires a genuinely
new method. The conjecture is verified for families on at most $12$ ground
elements~\cite{zivkovic2024} and for several restricted classes, but the
general case is open.

\subsection*{Contribution}

This paper records rigorous partial and structural results, each with a
complete proof:

\begin{enumerate}
\item Lemma~\ref{lem:singleton}: if $\{x\}\in\cF$ then $x$ is popular
  (density $\ge \Nhalf$).
\item Lemma~\ref{lem:injection}: for every $S\in\cF$,
  $|\{A:S\subseteq A\}|\ge|\{A:A\cap S=\emp\}|$, via an explicit injection.
\item Corollary~\ref{cor:degree}: $M\ge m/(k+1)$ where $M=\max_x d(x)$ and
  $k$ is the size of a smallest non-empty member-set.
\end{enumerate}

Together these show that any counterexample must have no singletons,
$k\ge2$, $\bar s<n/2$, and $M<m/2$. The jump from $M\ge m/(k+1)$ (below
$m/2$ whenever $k\ge2$) to the conjectured $M\ge m/2$ is precisely the open
gap.

We also report exact finite verification (computation as assistant, not proof)
for all union-closed families on $n\le 5$ ground elements; the tight family
is $\{\varnothing,[n]\}$, where every element has density exactly $\Nhalf$.

No claim is made that Frankl's conjecture is resolved, and none of the
results herein should be read as such. Whether Lemma~\ref{lem:singleton} or
Lemma~\ref{lem:injection} are novel is left to the literature; they are
presented as self-contained, verified proofs.

\section{Preliminaries}

\begin{definition}
For a union-closed family $\cF\subseteq\cP([n])$ with $m=|\cF|$ and
$x\in[n]$, define the \emph{frequency} (degree)
$$d(x):=|\{A\in\cF:x\in A\}|,\qquad
\widetilde d(x):=\frac{d(x)}{m}.$$
An element is \emph{popular} if $\widetilde d(x)\ge\Nhalf$.
For $S\subseteq[n]$ let
$$g(S):=|\{A\in\cF:S\subseteq A\}|,\qquad
h(S):=|\{A\in\cF:A\cap S=\emp\}|.$$
\end{definition}

\begin{observation}
For each $x$, the subfamily $\cF_x=\{A\in\cF:x\in A\}$ is itself
union-closed: if $A,B\in\cF_x$ then $A\cup B\in\cF$ and $x\in A\cup B$, so
$A\cup B\in\cF_x$. Likewise $\cF_{\bar x}=\{A\in\cF:x\notin A\}$ is
union-closed.
\end{observation}

\section{Structural results}

\begin{lemma}[Singleton case]
\label{lem:singleton}
If $\{x\}\in\cF$ then $d(x)\ge m/2$.
\end{lemma}

\begin{proof}
Define $\phi:\cF_{\bar x}\to\cF_x$ by $\phi(B)=B\cup\{x\}$.
It is well-defined because $B,\{x\}\in\cF$ and $\cF$ is union-closed, so
$B\cup\{x\}\in\cF$ and contains $x$. It is injective: if
$B_1\cup\{x\}=B_2\cup\{x\}$ then, since $x\notin B_1,B_2$, we get $B_1=B_2$.
Thus $|\cF_x|\ge|\cF_{\bar x}|$. As $\cF_x,\cF_{\bar x}$ partition $\cF$,
$m=d(x)+(m-d(x))$, so
$$d(x)\ge m-d(x)\iff 2d(x)\ge m\iff d(x)\ge\frac m2.$$
\end{proof}

\begin{lemma}[Injection bound]
\label{lem:injection}
For every $S\in\cF$, $g(S)\ge h(S)$.
\end{lemma}

\begin{proof}
Define $\phi:\{B\in\cF:B\cap S=\emp\}\to\{A\in\cF:S\subseteq A\}$ by
$\phi(B)=B\cup S$. Well-defined since $B,S\in\cF$ implies $B\cup S\in\cF$ and
$S\subseteq B\cup S$. Injective: if $B_1\cup S=B_2\cup S$ then, because
$B_1\cap S=B_2\cap S=\emp$, intersecting with $[n]\setminus S$ gives
$B_1=B_2$. Hence $g(S)\ge h(S)$.
\end{proof}

\begin{corollary}[Degree bound via smallest member-set]
\label{cor:degree}
Let $k=\min\{|S|:\emp\ne S\in\cF\}$ and $M=\max_x d(x)$. Then
$$M\ge \frac{m}{k+1}.$$
\end{corollary}

\begin{proof}
Choose $S\in\cF$ with $|S|=k$.
Every set containing $S$ contains some element of $S$, so
$\{A:S\subseteq A\}\subseteq\bigcup_{x\in S}\{A:x\in A\}$ and
$g(S)\le\sum_{x\in S}d(x)\le kM$; in particular $g(S)\le M$.
The sets meeting $S$ have cardinality at most $\sum_{x\in S}d(x)\le kM$, so
$$h(S)=m-\bigl|\{A:A\cap S\neq\emp\}\bigr|\ge m-kM.$$
By Lemma~\ref{lem:injection}, $M\ge g(S)\ge h(S)\ge m-kM$, hence
$(k+1)M\ge m$, i.e. $M\ge m/(k+1)$.
\end{proof}

\begin{corollary}
Any counterexample to Conjecture~\ref{conj:frankl} must satisfy:
\begin{enumerate}
\item no singleton member-sets ($k\ge2$, by Lemma~\ref{lem:singleton});
\item $\bar s<n/2$ (otherwise the average argument gives $M\ge m/2$);
\item $M<m/2$ (by definition of counterexample).
\end{enumerate}
If the counterexample has a 2-element member-set then $M\ge m/3$.
\end{corollary}

\section{Finite computation}

We verified Conjecture~\ref{conj:frankl} exactly for all non-trivial
union-closed families on $n=4$ and $n=5$ ground elements.

A union-closed family is determined by its inclusion-minimal non-empty
member-sets, an antichain $\mathcal M$: the family is $\{\varnothing\}$
together with all unions of subfamilies of $\mathcal M$. (Adding the empty
set only lowers every element's density fraction $d(x)/m$, so it is the
hardest case; families not containing $\varnothing$ are easier.) Enumerating
antichains therefore enumerates every union-closed family containing
$\varnothing$, and checking $2\max_x d(x)\ge m$ for each verifies the
conjecture for that ground size.

The computation uses exact integer arithmetic (bitmasks) and found no
counterexample:
\begin{center}
\begin{tabular}{c|c}
Ground set $n$ & Families checked \\ \hline
$4$ & $166$ \\
$5$ & $7\,579$
\end{tabular}
\end{center}
In every case the minimum of $\max_x d(x)/m$ is exactly $\Nhalf$, attained by
$\{\varnothing,[n]\}$, where each element has density exactly $\Nhalf$. This
shows the constant $\Nhalf$ in the conjecture is best possible.

This is computation as an \emph{assistant}: it verifies finite domains and
builds intuition, but proves nothing for general $n$. (The literature verifies
the conjecture up to $n=12$~\cite{zivkovic2024}.)

\section{Discussion and open questions}

The results above are partial/structural; the full conjecture remains open.
The strongest general bound known is $c\approx0.3823455$, and the entropy
method is sharp there, so new mathematics is needed to reach $\Nhalf$.

\begin{enumerate}
\item Can the bound $M\ge m/(k+1)$ be improved to $M\ge m/2$ for the case
  $k=2$ (families with a 2-element member-set) by a short argument? This is
  known in the literature but a self-contained proof would be worthwhile.
\item Is there a new method, outside the entropy framework, that closes the
  gap from $c\approx0.3823455$ to $\Nhalf$?
\item Extend the finite verification beyond $n=12$ or find a structural
  reduction that handles large $n$ from small cases.
\end{enumerate}

\section*{Acknowledgements and reproducibility}

All source code, proofs, and the research log are in the accompanying GitHub
repository (see README). Computations use only the Python standard library
(exact integer arithmetic); no proprietary software is required.

\begin{thebibliography}{9}
\bibitem{gilmer2022}
J.~Gilmer, ``A constant lower bound for Frankl's union-closed sets
conjecture,'' 2022.

\bibitem{sawin2022}
S.~Sawin, ``An improved lower bound for the union-closed set conjecture,''
arXiv:2212.12500, 2022.

\bibitem{liu2023}
J.~Liu, ``Improving the lower bound for the union-closed sets conjecture via
conditionally IID coupling,'' arXiv:2306.08824, 2023.

\bibitem{luraz2024}
K.~Lu and A.~Raz, ``Note on the union-closed sets conjecture and Reimer's
average set size theorem,'' arXiv:2405.10639, 2024.

\bibitem{cambie2025}
S.~Cambie, ``Better bounds for the union-closed sets conjecture using the
entropy approach,'' arXiv:2212.12500v2, 2025.

\bibitem{zivkovic2024}
M.~Zivkovi\'c and~?, ``The union-closed sets conjecture holds for families on
at most $12$ elements'' (finite verification).

\bibitem{reimer2003}
T.~Reimer, ``The average set size theorem,'' related to Reimer's theorem on
union-closed families.
\end{thebibliography}

\end{document}